% 93_repro.tex: Reproducibility Statement (fragment, no preamble)
% ICLR convention: unnumbered section placed before the references.
% This is the only place in the paper where repository-style artifact descriptions appear;
% categories only, no concrete file paths.

\section*{Reproducibility Statement}

All quantitative results in this paper were produced by a fixed set of pipeline and evaluation scripts, and every reported figure derives from archived per-sample JSON records: selection outputs, keyframe timestamps, grounding boxes, generated images, automatic metric values, and human labels are stored per sample, so each table can be recomputed from the raw records rather than from intermediate summaries, and the headline numbers were independently re-derived from these records in multiple verification passes. The main pipeline is training-free, and all reported runs were executed on a single RTX 4090D GPU with 24~GB of memory, which bounds the compute required for replication. The evaluation conventions that affect comparability, including the strict temporal pad of 0~s, the joint temporal-and-spatial metric, the no-box-counts-as-incorrect rule, and the bridging protocol for the external baseline, are stated where the corresponding results are reported (Section~\ref{sec:exp}). The human evaluation protocol, covering the faithfulness definition, the instruction tiers, and the annotation procedure, is archived together with the per-item labels. Clip identifiers for every evaluation subset, including the 282-clip common subset used in the selector comparison, are recorded so that all subsets can be reconstructed from the public V-STaR and HC-STVG~v2 releases. We will release the code (frame selection, grounding, generation, and evaluation scripts), the configuration files and prompts, and the per-sample records together with the paper.
